\section{Key-Recovery Security of Lyubashevsky Signatures} \label{sec-dilithium}

\adam{Ojaswi, Carter, Xinyu, Alex: This section merges the instantiability program with the circular-SIS working draft. The scheme below is the uncompressed warmup scheme of that draft with rejection sampling added, over $q=2^k$ with the gadget challenge set. Open items for the group are flagged inline.}

This part of the paper instantiates the hash function in a lattice-based Fiat--Shamir signature scheme, and proves the first two results of the three-step development of Part~I: key recovery from lossiness (this section), and static-message unforgeability from a circular assumption with obfuscation (Section~\ref{sec-dil-static}). The scheme is a Fiat--Shamir signature with aborts in the sense of Lyubashevsky~\cite{AC:Lyubashevsky09,EC:Lyubashevsky12}, the paradigm that Dilithium~\cite{Dilithium} instantiates and that motivates this part.

What transfers from this part to Dilithium proper is the instantiation technique, not the theorem statements. To make the difference precise we state at the outset what our scheme shares with Dilithium and what it does not. Shared: the Fiat--Shamir-with-aborts structure, an LWE-format public key, a masked commitment, and a response that verification checks to be short. Not shared: we work with plain matrices over $\Zq$ rather than over the ring $\Z_q[X]/(X^n+1)$, so there is no module structure, which is where Dilithium's efficiency comes from; our challenges are prefix-one binary decompositions rather than the sparse set $B_\tau$; our modulus is any sufficiently large prime rather than Dilithium's fixed NTT-friendly one; our signer is randomized, whereas ML-DSA specifies a deterministic mode and a hedged default; and we use neither compression of the commitment nor a hint vector. Section~\ref{sec-conclusion} records each of these differences as an open problem. The key-recovery result (Theorem~\ref{thm-dil-bk}) mirrors Section~\ref{sec-lf-kr}: the hash is a tunable lossy function composed with an outer hash, the proof switches to lossy mode and simulates signing by rejection sampling over the enumerated image, and the simulation is exact. It uses no obfuscation, and with a lattice-based lossy function it would be post-quantum throughout; constructing one with the required enumerability is an open problem we record in Section~\ref{sec-conclusion}. The static result (Theorem~\ref{thm-dil}) mirrors Section~\ref{sec-io-uf}. The hash function is an obfuscated circuit computing a masked conversion function, a forgery becomes a solution to the circular problem, and signing queries are handled by hardwiring. Two things change. First, the circular problem is $\CSIS$ with bit decomposition as the conversion function, and by Lemma~\ref{lem-csis-sis} this is not a new assumption; the reduction lands on plain SIS. Second, lattice verification and simulation include a low-norm slack term, which changes the exact steps of the Schnorr proof into steps with LWE and rejection-sampling accounting.

A standard-model proof also changes what must be modeled. Analyses of post-quantum schemes in idealized models must treat quantum access to the oracle~\cite{EC:KilLyuSch18}, and an instantiated hash function is a concrete circuit, so no oracle model remains in the statement. This does not by itself make the results post-quantum: the static construction uses iO, whose current instantiations are not post-quantum, and the key-recovery result would need both a quantum-secure lossy function and a proof stated against quantum adversaries. The static construction uses iO, whose current instantiations are not post-quantum, so that result is a feasibility statement about the shape of the scheme; the key-recovery result has no such caveat. The salted adaptive port is deferred to future work; the outlook at the end of Section~\ref{sec-dil-static} records the plan.

Throughout, $q$ is a prime, $k = \lceil \log_2 q\rceil$, and $\ell = 1 + nk$. The challenge space is $\chal_q$ and the conversion function is $\Fq_q$, both defined in Section~\ref{sec-lattice}; $\Fq_q$ is a bijection from $\Zq^n$ onto $\chal_q$, every challenge has infinity norm $1$, and no challenge is zero. Fix dimensions $n, \colsA$ and a mask bound $B$, and set the response bound $\Bclip := B - \ell$. We require $B \ge 4\ell(\colsA+n)$, which bounds the rejection probability of each signing attempt by $1/2$, and $q \ge B^2$, which keeps the modulus above every norm in the proofs. The key-recovery result works at any such $B$, including polynomial. The static result additionally requires the flooding condition of Section~\ref{sec-dil-static}, which requires $B$, and hence $q$, to be super-polynomial; quasi-polynomial suffices. \adam{Carter: Parameter sanity for your pass: for key recovery, $k = \Theta(\log\secp)$ closes the loop $\ell = nk$, $B \ge 4\ell(\colsA+n)$, $q = 2^k \ge B^2$; for the static result, $k = \Theta((\log\secp)^2)$ does. The final SIS norm is $B+\Bclip < 2B \le 2\sqrt{q}$. Concrete sizes need a real parameter table.}

\subsection{Lattice Problems} \label{sec-lattice}

This part works over $\Zq$ for a prime modulus $q$. Norms are the infinity norm, and $[\pm B]$ denotes the integer interval $\{-B,\ldots,B\}$. The following two problems are standard.

\begin{definition} \label{def-sis}
Let $q, n, m', \beta \in \N$. The advantage of an algorithm $\advA$ against $\SIS_{q,n,m',\beta}$ is
$$
\AdvSIS{q,n,m',\beta}{\advA}
:=
\Pr\big[\mat{A}\mat{x} = \mat{0} \,\land\, 0 < \|\mat{x}\|_\infty \le \beta
\;\big|\;
\mat{A} \getsr \Zq^{n \times m'},\ \mat{x} \getsr \advA(\mat{A})
\big] \;.
$$
\end{definition}

\begin{definition} \label{def-lwe}
Let $q, n, m', t, B \in \N$. The advantage of an algorithm $\advD$ against the decision problem $\LWE_{q,n,m',t,B}$ is
$$
\AdvLWE{q,n,m',t,B}{\advD}
:=
\big|
\Pr[\advD(\mat{A}, \mat{A}\mat{V}_1 + \mat{V}_2) \Rightarrow 1]
-
\Pr[\advD(\mat{A}, \mat{U}) \Rightarrow 1]
\big| \;,
$$
where $\mat{A} \getsr \Zq^{n \times m'}$, $\mat{V}_1 \getsr [\pm B]^{m' \times t}$, $\mat{V}_2 \getsr [\pm B]^{n \times t}$, and $\mat{U} \getsr \Zq^{n \times t}$.
\end{definition}

We also use the search version at the same parameters. The advantage of $\advA$ against search $\LWE_{q,n,m',t,B}$ is
$$
\AdvSLWE{q,n,m',t,B}{\advA}
:=
\Pr\big[\mat{A}\mat{V}_1' + \mat{V}_2' = \mat{T} \,\land\, \mat{V}_1' \in [\pm B]^{m'\times t} \,\land\, \mat{V}_2' \in [\pm B]^{n\times t}\big] \;,
$$
where $\mat{A} \getsr \Zq^{n\times m'}$, $\mat{V}_1 \getsr [\pm B]^{m'\times t}$, $\mat{V}_2 \getsr [\pm B]^{n\times t}$, $\mat{T} \gets \mat{A}\mat{V}_1+\mat{V}_2$, and $(\mat{V}_1',\mat{V}_2') \getsr \advA(\mat{A},\mat{T})$. The winner returns any consistent short pair, not necessarily the planted one.

\paragraph*{Discussion.} Definition~\ref{def-lwe} is decision LWE in normal form with $t$ independent secret--noise column pairs and noise uniform on a box. A hybrid over the columns reduces it to the single-column case with a factor $t$. SIS goes back to Ajtai~\cite{STOC:Ajtai96} and LWE to Regev~\cite{STOC:Regev05}. \adam{Ojaswi/Carter: Pin down the citation for LWE with uniform box noise and normal form at our parameter regimes; the working draft's Definition def:lwe is the same shape.}

The following circular problem is the SIS analogue of CDL. The conversion function output occurs both inside the equation and as the value the function takes on the solution point.

\begin{definition} \label{def-csis}
Let $q, n, m', \ell, \gamma, \eta \in \N$, let $\chal \subseteq \Z^{\ell}$, and let $f \colon \Zq^n \to \chal$ be efficient. The advantage of an algorithm $\advA$ against $\CSIS_{f,\chal_q,q,n,m',\ell,\gamma,\eta}$ is the probability that
$$
\mat{A}\hat{\mat{z}} - \mat{T}f(\hat{\mat{w}}) + \hat{\mat{u}} = \hat{\mat{w}},
\qquad
\|\hat{\mat{z}}\|_\infty \le \gamma,
\qquad
\|\hat{\mat{u}}\|_\infty \le \eta,
\qquad
(\hat{\mat{z}}, f(\hat{\mat{w}})) \ne \mat{0},
$$
where $\mat{A} \getsr \Zq^{n \times m'}$, $\mat{T} \getsr \Zq^{n \times \ell}$, and $(\hat{\mat{z}}, \hat{\mat{u}}, \hat{\mat{w}}) \getsr \advA(\mat{A}, \mat{T})$.
\end{definition}

\paragraph*{Discussion.} cSIS is a problem template; we write the leading c in lowercase to mark that it is a template, made concrete by a choice of $f$ and $\chal$. A choice of challenge space $\chal$ and conversion function $f$ makes it concrete, exactly as a choice of $f$ makes CDL concrete. The coordinate $\hat{\mat{u}}$ holds the low-norm slack of lattice verification. Schnorr verification is an exact group equation, so CDL has no such coordinate. Lattice verification includes a short error term, and that term appears in $\hat{\mat{u}}$. For the instantiation in this paper we use bit decomposition, and then cSIS is not a new assumption. It is SIS, by the following reduction.

We fix the conversion function used throughout this part. Let $q$ be any modulus and $k = \lceil\log_2 q\rceil$. For $a \in \Zq$ let $\binq_q(a) \in \bits^{k}$ be the binary representation of the integer representative of $a$ in $\{0,\ldots,q-1\}$, and let $D_q := \{\binq_q(a) : a \in \Zq\}$ be its image, a set of size $q$. Define
$$
\Fq_q(\mat{w}) := \big(1, \binq_q(w_1), \ldots, \binq_q(w_n)\big) \in \chal_q,
\qquad
\chal_q := \{1\} \times D_q^{\,n} \subseteq \bits^{\ell},
\qquad
\ell := 1 + nk,
$$
and let
$$
\mat{G}_q := \big[\, \mat{0}_{n\times 1} \;\big|\; \mat{I}_n \otimes (1, 2, \ldots, 2^{k-1}) \,\big] \in \Zq^{n\times\ell} \;.
$$
Four properties are immediate, and the proofs of this part use only these. The map $\Fq_q$ is a bijection from $\Zq^n$ onto $\chal_q$, so it takes the uniform distribution on $\Zq^n$ to the uniform distribution on $\chal_q$, for every modulus. We have $\mat{G}_q \Fq_q(\mat{w}) = \mat{w}$ for every $\mat{w}$, since the leading coordinate meets the zero column and the remaining coordinates reassemble each $w_i$ from its binary representation. Every element of $\chal_q$ has infinity norm $1$. And no element of $\chal_q$ is zero, because its first coordinate is $1$.

\paragraph*{Discussion.} The point of the leading coordinate and of taking the image of $\binq_q$ as the challenge space, rather than all of $\bits^{\ell}$, is that the proofs need a conversion function that takes a uniform masked point to a uniform challenge \emph{in the challenge space the scheme uses}. They do not need the challenge space to be every bit string of a given length. Binary decomposition is not balanced onto $\bits^{nk}$ at a prime modulus, but it is a bijection onto its own image at every modulus, which is what the argument uses. The leading $1$ then removes the zero challenge, which verification would otherwise have to reject and honest signing could otherwise produce.

\begin{lemma} \label{lem-csis-sis}
Let $q$ be prime and $f = \Fq_q$ with challenge space $\chal_q$ and $\ell = 1+nk$. Put $\alpha_{q,n} := \prod_{i=1}^{n}(1-q^{-i})$. For every algorithm $\advA$ there is an algorithm $\advB$ with the same running time up to an additive polynomial such that
$$
\AdvCSIS{\Fq_q,\chal_q,q,n,m',\ell,\gamma,\eta}{\advA}
\le
\alpha_{q,n}^{-1} \cdot \AdvSIS{q,n,m'+\ell+n,\max(\gamma,1,\eta)}{\advB} \;.
$$
\end{lemma}

\begin{proof}
Algorithm $\advB$ receives $\mat{A}' = [\mat{A} \mid \mat{B} \mid \mat{C}] \getsr \Zq^{n \times (m'+\ell+n)}$. If $\mat{C}$ is not invertible modulo $q$, it outputs $\bot$. Since $q$ is prime, a uniform $n\times n$ matrix over $\Zq$ is invertible with probability $\alpha_{q,n} = \prod_{i=1}^{n}(1-q^{-i})$, which is close to $1$ for a large prime. It sets $\mat{T} \gets -\mat{C}^{-1}\mat{B} - \mat{G}_q$ and $\mat{A}_0 \gets \mat{C}^{-1}\mat{A}$; conditioned on any invertible $\mat{C}$, the pair $(\mat{A}_0, \mat{T})$ is uniform. It runs $\advA(\mat{A}_0, \mat{T})$ and receives $(\hat{\mat{z}}, \hat{\mat{u}}, \hat{\mat{w}})$. Write $\mat{c} = \Fq_q(\hat{\mat{w}})$, so $\mat{G}_q\mat{c} = \hat{\mat{w}}$. The cSIS equation gives $\mat{A}_0\hat{\mat{z}} - \mat{T}\mat{c} + \hat{\mat{u}} = \mat{G}_q\mat{c}$, hence $\mat{A}_0\hat{\mat{z}} - (\mat{T} + \mat{G}_q)\mat{c} + \hat{\mat{u}} = \mat{0}$. Multiplying by $\mat{C}$ gives
$$
\mat{A}\hat{\mat{z}} + \mat{B}\mat{c} + \mat{C}\hat{\mat{u}} = \mat{0} \;,
$$
so $\advB$ returns $(\hat{\mat{z}}, \mat{c}, \hat{\mat{u}})$. Its entries are bounded by $\max(\gamma, 1, \eta)$, and it is nonzero because the first coordinate of $\mat{c}$ is $1$. Thus $\advB$ wins whenever $\mat{C}$ is invertible and $\advA$ wins. The instance $\advA$ receives is uniform for every invertible $\mat{C}$, so the two events are independent, and the claimed inequality follows from the invertibility probability.
\end{proof}

\subsection{The Scheme} \label{sec-dil-scheme}

A \emph{keyed hash family for $\Zq^n$} and message space $\bits^{\msglen}$ is a pair $\HF = (\HFKg, \HFEv)$ where $\HFKg$ on input a public key $(\mat{A}, \mat{T})$ outputs a key $\hk$, and the deterministic $\HFEv$ on inputs $\hk$ and $(\mat{w}, m) \in \Zq^n \times \bits^{\msglen}$ outputs an element of $\chal$. As in Section~\ref{sec-schnorr}, the hash key is generated honestly at key-generation time, may depend on the public key, and is placed in the verification key.

\begin{construction} \label{constr-dil}
Let $\HF$ be a keyed hash family for $\Zq^n$. The scheme $\DSScheme[\HF] = (\SchKeys, \SchSign, \SchVf)$ with message space $\bits^{\msglen}$ works as follows.

Algorithm $\SchKeys$ samples $\mat{A} \getsr \Zq^{n\times\colsA}$, $\mat{S}_1 \getsr [\pm1]^{\colsA\times\ell}$, $\mat{S}_2 \getsr [\pm1]^{n\times\ell}$, sets $\mat{T} \gets \mat{A}\mat{S}_1 + \mat{S}_2$, samples $\hk \getsr \HFKg((\mat{A},\mat{T}))$, and returns $\vkey = (\mat{A}, \mat{T}, \hk)$ and $\skey = (\mat{S}_1, \mat{S}_2, \vkey)$.

Algorithm $\SchSign$ on input $\skey, m$ repeats the following, at most $\secp$ times. Sample $\mat{y}_1 \getsr [\pm B]^{\colsA}$ and $\mat{y}_2 \getsr [\pm B]^{n}$, set
$$
\mat{w} \gets \mat{A}\mat{y}_1 + \mat{y}_2,
\qquad
c \gets \HFEv(\hk, (\mat{w}, m)),
\qquad
\mat{z}_i \gets \mat{y}_i + \mat{S}_i c \quad (i = 1,2) \;,
$$
and if $\|\mat{z}_1\|_\infty \le \Bclip$ and $\|\mat{z}_2\|_\infty \le \Bclip$, return $(c, \mat{z}_1, \mat{z}_2)$. If all $\secp$ attempts reject, return $\bot$.

Algorithm $\SchVf$ on inputs $\vkey, m, (c, \mat{z}_1, \mat{z}_2)$ computes $\mat{w} \gets \mat{A}\mat{z}_1 + \mat{z}_2 - \mat{T}c$ and returns 1 iff
$$
c = \HFEv(\hk, (\mat{w}, m)),
\qquad
\|\mat{z}_1\|_\infty \le \Bclip,
\qquad
\|\mat{z}_2\|_\infty \le \Bclip \;.
$$
\end{construction}

\paragraph*{Discussion.} This is the Fiat--Shamir transform of the Lyubashevsky-style identification protocol with an LWE-format key, uncompressed, with uniform-box masks and rejection sampling. No zero-challenge check is needed. The zero-free condition of Section~\ref{sec-io-uf} is met by construction here, since every element of $\chal_q$ has leading coordinate $1$ and is therefore nonzero. This is why verification checks only the hash equation and the two norm bounds. Rejection sampling with uniform boxes has an exactness property the proofs use. For any fixed value of $\mat{S}_ic$, the sum $\mat{y}_i + \mat{S}_ic$ with $\mat{y}_i$ uniform on $[\pm B]$, conditioned on landing in $[\pm\Bclip]$, is uniform on $[\pm\Bclip]$ and independent of $\mat{S}_ic$, because $\|\mat{S}_ic\|_\infty \le \ell$ shifts the box by at most $B - \Bclip$. For a fixed challenge, each attempt accepts with probability exactly $\rho := \left(\frac{2\Bclip+1}{2B+1}\right)^{\colsA+n}$, independent of the secret and the challenge. In honest signing the challenge depends on the attempt through the hash, and a weaker bound applies. Every $\mat{y}$ with $\|\mat{y}_i\|_\infty \le \Bclip-\ell$ is accepted whatever the challenge, so each attempt accepts with probability at least $1/2$ under $B \ge 4\ell(\colsA+n)$ (see also the discussion after Lemma~\ref{lem-dil-sim}), and signing returns $\bot$ with probability at most $2^{-\secp}$. Signing therefore meets the correctness requirement of Section~\ref{sec-prelims} with $\nu_{\mathsf{corr}} = 2^{-\secp}$, the probability that all $\secp$ attempts are rejected. No other failure mode remains: every challenge the signer computes lies in $\chal_q$, so verification never rejects an honestly produced signature for its challenge. Signing is randomized, which matches Dilithium's hedged variant; the other differences from Dilithium are listed at the start of Section~\ref{sec-dilithium}.

\subsection{The Construction} \label{sec-dil-bk}

The hash family composes the tunable lossy function with an outer hash into the challenge space. Let $\TLF$ be a tunable lossy function with domain $\bits^{\ellin}$, where $\ellin \ge nk + \msglen$ so that pairs $(\mat{w},m)$ embed by an injective encoding $\enc{\cdot}'' \colon \Zq^n \times \bits^{\msglen} \to \bits^{\ellin}$, and let $(\CRHFKg,\CRHFEv)$ be a collision-resistant keyed family from $\Ygrave$ to $\Zq^n$, as in Section~\ref{sec-lf-constr} but with range $\Zq^n$.

\begin{construction} \label{constr-dil-tlf}
Fix an image parameter $r\in[M]$. The keyed hash family $\HF^{\mathsf{ltlf}}_r=(\HFKg,\HFEv)$ for $\Zq^n$ and message space $\bits^{\msglen}$ is defined as follows. Algorithm $\HFKg$ on input $(\mat{A},\mat{T})$ computes $(K,\mathsl{td})\getsr\TLFKg(\usecp,r,\mathsf{inj})$ and $k'\getsr\CRHFKg(\usecp)$ and returns $\hk=(K,k')$; we write $k'$ because $k$ is the modulus exponent in this part. Algorithm $\HFEv$ on inputs $\hk$ and $(\mat{w},m)$ returns $\CRHFEv\bigl(k',\TLFEv(K,\enc{\mat{w},m}'')\bigr)$.
\end{construction}

\paragraph*{Discussion.} Key generation ignores $(\mat{A},\mat{T})$, and lossy keys occur only in the proof, exactly as in Section~\ref{sec-lf-kr}.

\subsection{Main Result} \label{sec-dil-bk-thm}

The following lemma is the lattice analogue of Lemma~\ref{lem-sim}, and it is exact. It contains the reason rejection sampling and lossy simulation compose: the secret is recoverable from an accepted response, so accepted attempts correspond one-to-one with the pairs the simulator samples. We write $\mat{y} = (\mat{y}_1, \mat{y}_2)$, $\mat{z} = (\mat{z}_1, \mat{z}_2)$, and $\mat{S}c := (\mat{S}_1c, \mat{S}_2c)$, so that $\mat{z} = \mat{y} + \mat{S}c$ componentwise.

\begin{lemma} \label{lem-dil-sim}
Fix $\mat{A}$, short $\mat{S}_1,\mat{S}_2$ with entries in $[\pm1]$, $\mat{T}=\mat{A}\mat{S}_1+\mat{S}_2$, a function $H \colon \Zq^n \times \bits^{\msglen} \to \chal_q$, a message $m$, and a finite set $C \subseteq \bits^{\ell}$ with $\Image(H) \subseteq C$. Call a pair $(c,\mat{z})$ with $\mat{z}=(\mat{z}_1,\mat{z}_2)$ \emph{good} if
$$
\|\mat{z}_1\|_\infty \le \Bclip,
\qquad
\|\mat{z}_2\|_\infty \le \Bclip,
\qquad
H\big(\mat{A}\mat{z}_1+\mat{z}_2-\mat{T}c,\ m\big)=c \;.
$$
Consider the honest attempt: $\mat{y} \getsr [\pm B]^{\colsA+n}$, $c \gets H(\mat{A}\mat{y}_1+\mat{y}_2, m)$, $\mat{z}_i \gets \mat{y}_i+\mat{S}_ic$, accepted if both norms are at most $\Bclip$. Consider the simulated attempt: $c \getsr C$, $\mat{z}$ uniform on the clipped boxes, accepted if $(c,\mat{z})$ is good. Then the map $\mat{y} \mapsto (c,\mat{z})$ is a bijection from accepted honest attempts to good pairs, and both attempts, conditioned on acceptance, output a uniformly distributed good pair.
\end{lemma}

\begin{proof}
The honest attempt lands in the good set: acceptance bounds the norms, and $\mat{A}\mat{z}_1+\mat{z}_2-\mat{T}c = \mat{A}\mat{y}_1+\mat{y}_2 + (\mat{A}\mat{S}_1+\mat{S}_2)c - \mat{T}c = \mat{A}\mat{y}_1+\mat{y}_2$, so the good-pair condition is the definition of $c$. The map is injective because $\mat{y} = \mat{z}-\mat{S}c$ recovers the attempt. It is surjective: given a good pair, set $\mat{y} \gets \mat{z}-\mat{S}c$; then $\|\mat{y}_i\|_\infty \le \Bclip + \ell \le B$, so $\mat{y}$ is in the attempt space, the commitment $\mat{A}\mat{y}_1+\mat{y}_2$ equals $\mat{A}\mat{z}_1+\mat{z}_2-\mat{T}c$ by the same computation, the good-pair condition makes the attempt's challenge equal $c$, and the attempt is accepted with response $\mat{z}$. Each accepted attempt arises from exactly one uniformly chosen $\mat{y}$, so the honest conditional distribution is uniform on good pairs. The simulated attempt is uniform on $C \times$ the clipped boxes, and its acceptance event is exactly the good set, because every good pair has $c \in \Image(H) \subseteq C$; conditioning gives the uniform distribution on the same set.
\end{proof}

\paragraph*{Discussion.} Two counting facts make the simulation efficient, and both are independent of $H$. First, every $\mat{y}$ with $\|\mat{y}_i\|_\infty \le \Bclip-\ell$ is accepted regardless of the challenge, since the shift $\mat{S}_ic$ moves it by at most $\ell$; hence the honest attempt accepts with probability at least $\big(\tfrac{2(\Bclip-\ell)+1}{2B+1}\big)^{\colsA+n} \ge 1/2$ under $B \ge 4\ell(\colsA+n)$, and the number of good pairs is at least half the number of clipped responses. Second, therefore, the simulated attempt accepts with probability at least $1/(2|C|)$, so $2\secp|C|$ trials fail with probability at most $e^{-\secp}$.

\begin{theorem} \label{thm-dil-bk}
Let $\advA$ be an adversary against BK-CMA of $\DSScheme[\HF^{\mathsf{ltlf}}_r]$ that runs in time $t_{\advA}$ and makes at most $q_s$ signing queries. For every $r\in[M]$, there are a distinguisher $\advD$ and an algorithm $\advB$ such that
$$
\AdvBKCMA{\DSScheme[\HF^{\mathsf{ltlf}}_r]}{\advA}
\le
\AdvTLF{\TLF,r}{\advD}
+\AdvSLWE{q,n,\colsA,\ell,1}{\advB}
+\nu_{\mathsf{enum}}(\secp,r)
+q_s\big(e^{-\secp}+2^{-\secp}\big) \;.
$$
Adversary $\advD$ runs in time $t_{\advA}+q_s\poly(\secp)$, and $\advB$ runs in time $t_{\advA}+\tau(\secp,r)+(q_s+1)\secp\, r\poly(\secp)$.
\end{theorem}

\begin{proof}
The games mirror Theorem~\ref{thm-bk}. Game $\Gm_0$ is BK-CMA. Game $\Gm_1$ generates the TLF key in lossy mode at the same $r$; the distinguisher $\advD$ samples $\mat{A},\mat{S}_1,\mat{S}_2$ and the outer-hash key itself, signs honestly, and outputs 1 iff $\advA$ returns the signing key, so the switch costs $\AdvTLF{\TLF,r}{\advD}$. Game $\Gm_2$ computes $S\getsr\TLFEnum(\mathsl{td})$ and $C\gets\{\CRHFEv(k,y):y\in S\}$ once, and answers each signing query by up to $2\secp|C|$ simulated attempts of Lemma~\ref{lem-dil-sim}, returning $\bot$ if all fail. In lossy mode every hash value lies in $C$, except with probability $\nu_{\mathsf{enum}}(\secp,r)$, so the lemma applies: conditioned on producing an output, each query is answered with exactly the distribution of $\Gm_1$. The output distributions differ only through the two failure events, honest signing exhausting $\secp$ attempts and simulation exhausting $2\secp|C|$ trials, which the counting facts above bound by $2^{-\secp}$ and $e^{-\secp}$ per query. Finally, algorithm $\advB$ receives a search-LWE challenge $(\mat{A},\mat{T})$ with $\ell$ columns at box $1$, which has exactly the distribution of the public key, samples the lossy hash key itself, answers signing queries as in $\Gm_2$ without any secret, and forwards the matrix pair from $\advA$'s output key. If $\advA$ wins BK-CMA, that pair is the planted $(\mat{S}_1,\mat{S}_2)$, which is a valid search-LWE solution. The reduction in fact needs less. It succeeds whenever $\advA$ outputs any pair $(\mat{S}_1',\mat{S}_2')$ of matrices with entries in $[\pm1]$ satisfying $\mat{A}\mat{S}_1'+\mat{S}_2' = \mat{T}$, whether or not it equals the planted one. Theorem~\ref{thm-dil-bk} therefore also holds for the stronger notion in which the adversary wins by producing any such pair, which is the more natural notion here, since the key is not determined by the public key. Hence $\Pr[\Gm_2\Rightarrow1]\le\AdvSLWE{q,n,\colsA,\ell,1}{\advB}$, and the running-time claims follow by inspection.
\end{proof}

\paragraph*{Discussion.} The interpretation matches Section~\ref{sec-lf-kr}: with this hash family, the signing oracle gives no help toward key recovery beyond the public key itself, at the price of assuming search LWE hard for time $r \cdot \mathrm{poly}$. The theorem is generic in the tunable lossy function. A lattice-based instantiation would make every ingredient lattice-based. The candidate is the LWE-based construction of Peikert and Waters at a fixed quasi-polynomial image bound, whose enumerability remains to be verified, and Section~\ref{sec-conclusion} records this, together with the universal-injective-mode question, as an open problem. With such an instantiation the assumption is search LWE against quasi-polynomial time. No obfuscation occurs anywhere in this section, and the reductions are straight-line, without rewinding, so we expect the result to hold against quantum adversaries that make classical signing queries. \adam{Ojaswi/Carter: Two items: (i) verify the counting-fact constants and the $2\secp|C|$ trial bound against the write-up of Theorem~\ref{thm-bk}; (ii) the BK-CMA win condition of Figure~\ref{fig-eufcma} is exact key equality, which is a special case of the search-LWE solution condition, so the last reduction step is sound as stated, but consider whether to strengthen the theorem to the relaxed win condition (any consistent short key), which the same proof gives.}

\section{Bounded Static-Message Unforgeability of Lyubashevsky Signatures} \label{sec-dil-static}

We now instantiate the hash function of the scheme of Construction~\ref{constr-dil} with an obfuscated circuit and prove static-message unforgeability, mirroring Section~\ref{sec-io-uf}. Throughout this section, fix a polynomial $Q = Q(\secp)$; as there, the scheme is parameterized by $Q$ through circuit padding, and the theorem covers adversaries declaring at most $Q$ signing messages. Throughout this section we additionally impose the \emph{flooding condition}: $\ell(\colsA+n)/B$ is negligible, for example $B = \ell(\colsA+n)\cdot 2^{(\log\secp)^2}$. Under it, a signing attempt rejects with negligible probability, so each query is answered by its first attempt except with negligible probability, and the point the proof hardwires is fixed by the presampled nonce alone, before any hash value exists. The Discussion at the end of Section~\ref{sec-dil-thm} explains why the proof needs this.

\subsection{The Construction} \label{sec-dil-constr}

Let $\PRF_1$ be a puncturable PRF with domain $\bits^{\ell_2}$ and range $[\pm B]^{\colsA}$, and let $\PRF_2$ be a puncturable PRF with domain $\bits^{\ell_2}$ and range $[\pm B]^{n}$, where $\enc{\cdot}''' \colon \Zq^n \times \bits^{\msglen} \to \bits^{\ell_2}$ is an injective encoding. We treat the PRF ranges as exactly uniform on the boxes; standard constructions give ranges statistically close to uniform, which adds a term of the form $q_s \cdot \mathrm{negl}(\secp)$ that we suppress. Let $\iO$ be an indistinguishability obfuscator.

\begin{figure}[t]
\begin{center}
\fbox{
\begin{minipage}{4.6in}
\begin{tabbing}
123\=123\=123\=\kill
\textbf{Circuit} $\CircD[\mat{A}, \mat{T}, k_1, k_2](\mat{w}, m)$: \\
\> $\mat{v}_1 \gets \PRFEv_1(k_1, \enc{\mat{w},m}''') \;;\; \mat{v}_2 \gets \PRFEv_2(k_2, \enc{\mat{w},m}''')$ \\
\> Return $\Fq_q\big(\mat{w} + \mat{A}\mat{v}_1 + \mat{v}_2\big)$
\end{tabbing}
\end{minipage}
}
\end{center}
\vspace{-6pt}
\caption{The lattice hash circuit underlying $\HF^{\mathsf{lio}}_Q$. It is padded to the size bound $s_{\mathrm{D}}(\secp, Q)$ before obfuscation.}\label{fig-circuit-dil}
\vspace{-3pt}\hrulefill
\vspace{-1ex}
\end{figure}

\begin{construction} \label{constr-dil-hash}
The keyed hash family $\HF^{\mathsf{lio}}_Q = (\HFKg, \HFEv)$ for $\Zq^n$ and message space $\bits^{\msglen}$ is defined as follows. Algorithm $\HFKg$ on input $(\mat{A}, \mat{T})$ computes $k_i \getsr \PRFKg_i(\usecp)$ for $i \in [2]$ and $\widetilde{\CircD} \getsr \iO(\usecp, \CircD[\mat{A}, \mat{T}, k_1, k_2])$ for the circuit $\CircD$ in Figure~\ref{fig-circuit-dil}, padded to size $s_{\mathrm{D}}(\secp, Q)$, and returns $\hk \gets \widetilde{\CircD}$. Algorithm $\HFEv$ evaluates the obfuscated circuit. Here $s_{\mathrm{D}}(\secp, Q)$ is the fixed polynomial bounding the size of the circuits in the proof of Theorem~\ref{thm-dil}.
\end{construction}

\paragraph*{Discussion.} The circuit is the masked-conversion simulator of the ROM proof made concrete, exactly as the circuit of Figure~\ref{fig-circuit} is for Schnorr, with one structural difference in where uniformity of hash values comes from. The Schnorr circuit outputs $f(R X^a g^b) + a$, and the added mask $a$ makes the output exactly uniform for any $f$. Here the challenge space consists of short vectors, and adding a uniform challenge-space element would destroy shortness, so no output mask is available. Uniformity comes instead from the conversion function itself. The masked point $\mat{w} + \mat{A}\mat{v}_1 + \mat{v}_2$ is pseudorandom under LWE, and $\Fq_q$ maps the uniform distribution on $\Zq^n$ to the uniform distribution on $\chal_q$, because it is a bijection onto $\chal_q$ at every modulus. This is what the prefix-one challenge space provides, and no balance condition on the modulus is needed.

The extraction identity is the analogue of the CDL conversion. If $(c, \mat{z}_1, \mat{z}_2)$ verifies for $m$, then $\mat{A}\mat{z}_1 + \mat{z}_2 - \mat{T}c = \mat{w}$ and $c = \Fq_q(\mat{w}')$ with $\mat{w}' = \mat{w} + \mat{A}\mat{v}_1 + \mat{v}_2$. Adding $\mat{A}\mat{v}_1 + \mat{v}_2$ to the verification equation gives
$$
\mat{A}(\mat{z}_1 + \mat{v}_1) - \mat{T}\,\Fq_q(\mat{w}') + (\mat{z}_2 + \mat{v}_2) = \mat{w}' \;,
$$
which is a $\CSIS$ solution $(\mat{z}_1{+}\mat{v}_1,\ \mat{z}_2{+}\mat{v}_2,\ \mat{w}')$ with norms at most $\Bclip + B$, nontrivial because the first coordinate of $c$ is $1$. The verifier's slack $\mat{z}_2$ and the mask $\mat{v}_2$ end up in the error coordinate of the solution; this is the reason $\CSIS$ has that coordinate. The reduction generates the PRF keys itself, so it can recompute $(\mat{v}_1, \mat{v}_2)$ from the forgery.

\subsection{Main Result} \label{sec-dil-thm}

\begin{theorem} \label{thm-dil}
Let $q$ be prime, let $\chal_q$ and $\Fq_q$ be as in Section~\ref{sec-lattice} with $\ell = 1+nk$, and let $f = \Fq_q$. Let $\PRF_1$ and $\PRF_2$ be puncturable PRFs with the domains and ranges of Section~\ref{sec-dil-constr}, and let $\iO$ be an indistinguishability obfuscator. Let $\advA = (\advA_0, \advA_1)$ be an adversary that runs in time $t_{\advA}$ and declares $q_s \le Q$ signing messages. There are an adversary $\advB$ against $\CSIS$, a distinguisher $\advD_{\mathsf{io}}$ against $\iO$, distinguishers $\advD^{\mathsf{prf}}_1, \advD^{\mathsf{prf}}_2$ for the PRFs, and distinguishers $\advD_1, \advD_2$ against decision LWE, each running in time $t_{\advA} + \poly(\secp, Q)$, such that
$$
\begin{aligned}
\AdvEUFSTATICQ{\DSScheme[\HF^{\mathsf{lio}}_Q]}{\advA}
\le\ &\AdvCSIS{\Fq_q,\chal_q,q,n,\colsA,\ell,\Bclip+B,\Bclip+B}{\advB}
+ \AdvIO{}{\advD_{\mathsf{io}}} \\
&+ \sum_{i=1}^{2}\AdvPRF{\PRF_i}{\advD^{\mathsf{prf}}_i}
+ \AdvLWE{q,n,\colsA,q_s,B}{\advD_1}
+ \AdvLWE{q,n,\colsA,\ell,1}{\advD_2} \\
&+ q_s \cdot \frac{4\ell(\colsA+n)}{2B+1}
+ q_s^2(2B+1)^{-n} \;.
\end{aligned}
$$
\end{theorem}

\begin{corollary} \label{cor-dil-sis}
Assume the setting of Theorem~\ref{thm-dil}, and suppose $\ell(\colsA+n)/B$ is negligible, $\SIS_{q,n,\colsA+\ell+n,\Bclip+B}$ holds, decision LWE holds at the parameters of the theorem, and the PRFs and $\iO$ are secure. Then $\DSScheme[\HF^{\mathsf{lio}}_Q]$ is $\mathrm{EUF\text{-}static}[Q]$ secure, for each fixed polynomial $Q$; as in Section~\ref{sec-io-uf}, this is a statement about the family indexed by $Q$ rather than ordinary static security of one scheme.
\end{corollary}

The corollary follows from Theorem~\ref{thm-dil} and Lemma~\ref{lem-csis-sis}; the factor $4$ of the lemma is absorbed asymptotically. The assumption profile is SIS, LWE, puncturable PRFs, and the assumptions underlying iO, all polynomially secure, at a quasi-polynomial modulus, since the flooding condition requires $B$, and hence $q \ge B^2$, to be super-polynomial, and quasi-polynomial suffices. To our knowledge this is the first standard-model static-message unforgeability result for a lattice-based Fiat--Shamir signature scheme.

\begin{proof}[of Theorem~\ref{thm-dil}]
The proof follows the game chain of Theorem~\ref{thm-uf}, and we describe each game by its difference from that proof, writing $m_1, \ldots, m_{q_s}$ for the declared messages with $q_s \le Q$ and $m^*$ for the declared forgery message.

\heading{Game $\Gm_0$ (the static game).} The game presamples the $q_s$ signatures before building the circuit, running the signing algorithm with direct PRF evaluation in place of the obfuscated circuit; this is the same computation, so nothing changes.

\heading{Game $\Gm_1$ (rejection abort).} Abort, with output 0, if the first attempt of any signing query rejects. A coordinate of $\mat{z}_i = \mat{y}_i + \mat{S}_ic$ can exceed $\Bclip$ only if the corresponding coordinate of $\mat{y}_i$ exceeds $B - 2\ell$ in absolute value, whatever the challenge is, so one attempt rejects with probability at most $4\ell(\colsA+n)/(2B+1)$, and the abort costs at most $q_s \cdot 4\ell(\colsA+n)/(2B+1)$, which the flooding condition makes negligible. On non-aborting executions every query is answered by its first attempt. Write $(\mat{w}_i, c_i, \mat{z}_i)$ for that attempt; its point $(\mat{w}_i, m_i)$ is a function of the presampled $\mat{y}_i$ alone, and in particular is independent of every hash value. This independence is what the flooding condition provides, and the hops below rely on it.

\heading{Game $\Gm_2$ (collision abort).} Abort, with output 0, if two of the $q_s$ points $(\mat{w}_i, m_i)$ collide. Each $\mat{w}_i = \mat{A}\mat{y}_{i,1} + \mat{y}_{i,2}$ has, conditioned on everything sampled independently of $\mat{y}_{i,2}$, min-entropy at least that of $\mat{y}_{i,2}$, so any fixed value is hit with probability at most $(2B+1)^{-n}$. A union over at most $q_s^2$ pairs bounds the abort by $q_s^2 (2B+1)^{-n}$.

\heading{Game $\Gm_3$ (hardwiring).} Replace the obfuscated circuit by an obfuscation of the circuit $\CircD'$ that contains the table $\{(\mat{w}_i, m_i) \mapsto c_i\}_{i \in [q_s]}$, punctured keys $k_1\{\mathcal{P}\}, k_2\{\mathcal{P}\}$ for the puncture set $\mathcal{P} = \{\enc{\mat{w}_i, m_i}''\}$, and the main branch of $\CircD$ otherwise, padded to $s_{\mathrm{D}}(\secp, Q)$. With the true values hardwired the two circuits compute the same function, exactly as in Game~$\Gm_2$ of Theorem~\ref{thm-uf}, so iO bounds the change.

\heading{Game $\Gm_4$ (uniform masks).} Replace the outputs of $\PRF_1$ and then $\PRF_2$ at the $q_s$ points of $\mathcal{P}$ by independent uniform values on their ranges. Two puncturable-PRF hops, as in Game~$\Gm_3$ of Theorem~\ref{thm-uf}. The puncture set is fixed by the presampled nonces, before any PRF challenge value is used, so the selective puncturing requirement of Definition~\ref{def-pprf} is met.

\heading{Game $\Gm_5$ (uniform hash values).} For each $i$, the hardwired value is $c_i = \Fq_q(\mat{w}_i + \mat{A}\mat{v}_{1,i} + \mat{v}_{2,i})$ with $(\mat{v}_{1,i}, \mat{v}_{2,i})$ uniform on boxes and used nowhere else. First replace the $q_s$ masked points $\mat{w}_i + \mat{A}\mat{v}_{1,i} + \mat{v}_{2,i}$ by independent uniform elements of $\Zq^n$; this is decision LWE with matrix $\mat{A}$, $q_s$ secret--noise pairs, and box $B$, at a cost of $\AdvLWE{q,n,\colsA,q_s,B}{\advD_1}$. Then $c_i = \Fq_q(\text{uniform})$ is uniform on $\chal_q$, exactly, because $\Fq_q$ is a bijection onto $\chal_q$; discard the masks. This is the analogue of the exact step of Game~$\Gm_4$ of Theorem~\ref{thm-uf}, with the exactness split into one LWE hop and one exact bijection step. The masked points are shifted by the $\mat{w}_i$, which are independent of the masks, so the shift preserves the LWE reduction.

\heading{Game $\Gm_6$ (signing without the secret).} For each $i$, the record is now produced as follows. With probability $1-\rho$, where $\rho = \big(\tfrac{2\Bclip+1}{2B+1}\big)^{\colsA+n}$ is the fixed-challenge acceptance probability of Section~\ref{sec-dil-scheme}, trigger the rejection abort of $\Gm_1$. Otherwise sample $c_i \getsr \chal_q$ and $\mat{z}_i$ uniform on $[\pm\Bclip]^{\colsA} \times [\pm\Bclip]^{n}$, set $\mat{w}_i \gets \mat{A}\mat{z}_{i,1} + \mat{z}_{i,2} - \mat{T}c_i$, hardwire $(\mat{w}_i, m_i) \mapsto c_i$, and return $(c_i, \mat{z}_i)$; the collision abort of $\Gm_2$ is evaluated on these re-derived $\mat{w}_i$. This has the same distribution as $\Gm_5$, exactly. In $\Gm_5$ the value $c_i$ is uniform and independent of $\mat{y}_i$, so the first attempt rejects with probability exactly $1-\rho$, independent of the secret and the challenge, and, conditioned on acceptance, $\mat{z}_i = \mat{y}_i + \mat{S}_ic_i$ is uniform on the clipped box and independent of $\mat{S}_ic_i$, by the rejection-sampling exactness of Section~\ref{sec-dil-scheme}; the point $\mat{w}_i$ is then determined by the identity $\mat{w}_i = \mat{A}\mat{z}_{i,1} + \mat{z}_{i,2} - \mat{T}c_i$. This is the analogue of the change of variables of Game~$\Gm_5$ of Theorem~\ref{thm-uf}; there the identity is $R_i = g^{s_i}X^{-v_i}$, here it is the verification equation. The secrets $\mat{S}_1, \mat{S}_2$ are no longer used in signing.

\heading{Game $\Gm_7$ (uniform public key).} Replace $\mat{T} = \mat{A}\mat{S}_1 + \mat{S}_2$ by uniform $\mat{T} \getsr \Zq^{n\times\ell}$. Signing no longer uses the secrets, so this is decision LWE with $\ell$ secret--noise columns at box $1$, at a cost of $\AdvLWE{q,n,\colsA,\ell,1}{\advD_2}$.

\heading{Forgery extraction.} In $\Gm_7$ the pair $(\mat{A}, \mat{T})$ is uniform, which is exactly a $\CSIS$ instance. The reduction embeds it, simulates $\Gm_7$, and receives a valid forgery $(c^*, \mat{z}^*)$ on $m^*$ with probability $\Pr[\Gm_7 \Rightarrow 1]$. The point $(\mat{w}^*, m^*)$ with $\mat{w}^* = \mat{A}\mat{z}^*_1 + \mat{z}^*_2 - \mat{T}c^*$ is not in the hardwired table: the message $m^*$ is fresh, and the table is indexed by pairs. So the circuit computes $c^*$ through the main branch, and the extraction identity of Section~\ref{sec-dil-constr} gives the $\CSIS$ solution $(\mat{z}^*_1{+}\mat{v}^*_1,\ \mat{z}^*_2{+}\mat{v}^*_2,\ \mat{w}'^*)$ with norms at most $\Bclip + B$, nontrivial since verification rejects $c^* = 0^\ell$. Hence $\Pr[\Gm_7 \Rightarrow 1] \le \AdvCSIS{\Fq_q,\chal_q,q,n,\colsA,\ell,\Bclip+B,\Bclip+B}{\advB}$.

\heading{Conclusion.} Collecting,
$$
\begin{aligned}
\AdvEUFSTATICQ{\DSScheme[\HF^{\mathsf{lio}}_Q]}{\advA}
\le\ &\AdvCSIS{\Fq_q,\chal_q,q,n,\colsA,\ell,\Bclip+B,\Bclip+B}{\advB}
+ \AdvIO{}{\advD_{\mathsf{io}}} \\
&+ \sum_{i=1}^{2}\AdvPRF{\PRF_i}{\advD^{\mathsf{prf}}_i}
+ \AdvLWE{q,n,\colsA,q_s,B}{\advD_1}
+ \AdvLWE{q,n,\colsA,\ell,1}{\advD_2} \\
&+ q_s \cdot \frac{4\ell(\colsA+n)}{2B+1}
+ q_s^2(2B+1)^{-n} \;,
\end{aligned}
$$
where the single iO term is the hardwiring switch of $\Gm_3$; the later hops change only the distribution of the hardwired data, which their own reductions build and obfuscate themselves, so no further iO invocation occurs. This is the bound of the theorem, and the running-time claims follow by inspection.
\end{proof}

\paragraph*{Discussion.} The flooding condition is what fixes the hardwired points before any hash value is computed. With rejection at polynomial $B$, the identity of the accepted attempt depends on the hash values, which are what the hybrids replace, so the puncture set is not fixed in advance and the selective puncturing argument does not apply. Moreover the rejected attempts' pairs are correlated with the secret, which blocks the key switch. We do not know a proof in that regime, and we record it as an open problem in Section~\ref{sec-conclusion}. The key-recovery result has no such condition, because its simulation never punctures. \adam{All: Internal review note: the polynomial-$B$ static claim was withdrawn after we found the acceptance/puncturing circularity and the secret-correlation of rejected attempts; do not re-attempt it casually. The flooding regime keeps all primitives polynomially secure; the cost is a quasi-polynomial modulus.}

Three further remarks, mirroring Section~\ref{sec-io-uf}. The parameter $Q$ enters only through the padding bound $s_{\mathrm{D}}(\secp, Q)$, and signing is stateless and unrestricted. The reduction is non-black-box in the hash function, using the PRF keys to recompute the masks of the forgery, which is how the result relates to the impossibility landscape for Fiat--Shamir in the standard model. The construction is a feasibility result; evaluating an obfuscated circuit per hash call is not a practical proposal.

\subsection{Removing the Query Bound: the Derandomized Scheme} \label{sec-dil-derand}

\adam{All: New material, written by one hand and unread by anyone else, like Section~\ref{sec-derand} was. The five spots that most want independent checking are listed in the note after the proof.}

This subsection ports the derandomization of Section~\ref{sec-derand} to the lattice scheme. The signer derives its nonces from the signing key and the message, which is how ML-DSA's deterministic mode already signs, and the hash circuit gets a trigger branch indexed by the message. The hardwired table of Theorem~\ref{thm-dil} disappears, and with it the parameter $Q$: the result is EUF-static security with no query bound, under the same assumptions and the same flooding condition.

\paragraph*{The scheme.} Let $\PRF_0$ be a pseudorandom function with range $[\pm B]^{\colsA} \times [\pm B]^{n}$, applied to an injective encoding of $(m, \kappa)$ for an attempt counter $\kappa \in [\secp]$; it need not be puncturable. The scheme $\DSScheme^{\mathsf{d}}[\HF]$ is the scheme of Construction~\ref{constr-dil} with two changes. Key generation additionally samples $k_0 \getsr \PRFKg_0(\usecp)$ and stores it in the signing key. Attempt $\kappa$ of signing computes its nonces as $(\mat{y}_1, \mat{y}_2) \gets \PRFEv_0(k_0, \enc{m,\kappa})$ instead of sampling them. Verification is unchanged. Signing is deterministic, so a repeated query is answered by replaying, exactly as in Section~\ref{sec-derand}, and ML-DSA's deterministic mode runs the same counter loop.

\paragraph*{The hash construction.} Let $\PRF_3$ be a puncturable PRF with domain $\bits^{\msglen}$ and range $\Zq^n$, and let $\PRF_4$ be a puncturable PRF with domain $\bits^{\msglen}$ and range $[\pm\Bclip]^{\colsA} \times [\pm\Bclip]^{n}$. The circuit $\CircD^{\mathsf{d}}[\mat{A}, \mat{T}, k_1, k_2, k_3, k_4](\mat{w}, m)$ computes
$$
c_t \gets \Fq_q\big(\PRFEv_3(k_3, m)\big),
\qquad
(\mat{z}_{t,1}, \mat{z}_{t,2}) \gets \PRFEv_4(k_4, m),
$$
returns $c_t$ if $\mat{w} = \mat{A}\mat{z}_{t,1} + \mat{z}_{t,2} - \mat{T}c_t$, and otherwise returns $\Fq_q(\mat{w} + \mat{A}\mat{v}_1 + \mat{v}_2)$ with $\mat{v}_1, \mat{v}_2$ computed from $k_1, k_2$ as in Figure~\ref{fig-circuit-dil}. The keyed hash family $\HF^{\mathsf{dlio}}$ obfuscates this circuit, padded to a fixed polynomial size bound $s^{\mathsf{d}}_{\mathrm{D}}(\secp)$ that does not depend on any query count, because the circuits in the proof below hardwire two points, whatever the number of declared messages. Since $\Fq_q$ is a bijection onto $\chal_q$, the trigger challenge $c_t$ is uniform on $\chal_q$ when $\PRFEv_3(k_3,m)$ is uniform on $\Zq^n$.

\paragraph*{Discussion.} For a slot $m$, the pair $(c_t, \mat{z}_t)$ with $\mat{W}_m := \mat{A}\mat{z}_{t,1} + \mat{z}_{t,2} - \mat{T}c_t$ is a valid signature on $m$, computed from the public key and the PRF keys alone. This is how the reduction signs. An honest signing attempt reaches the trigger point with probability at most $(2B+1)^{-n}$, because $\mat{y}_2$ gives the commitment that much min-entropy at any fixed target, and a signature produced through the trigger branch verifies in any case.

\begin{theorem} \label{thm-dil-der}
Impose the flooding condition. Let $\advA = (\advA_0, \advA_1)$ be an adversary that runs in time $t_{\advA}$ and declares $q_s$ signing messages, where $q_s$ is not bounded by any parameter of the scheme. There are adversaries and distinguishers, each running in time $t_{\advA} + \poly(\secp, q_s)$, such that, writing $\alpha_{q,n} = \prod_{i=1}^{n}(1-q^{-i})$ and
$$
\epsilon_j
:=
\sum_{\beta\in\{0,1\}}
\Big(
\AdvIO{}{\advD^{\mathsf{io}}_{j,\beta}}
+\sum_{i=1}^{4}\AdvPRF{\PRF_i}{\advD^{\mathsf{prf}}_{j,\beta,i}}
+\AdvLWE{q,n,\colsA,1,B}{\advD^{\mathsf{lwe}}_{j,\beta}}
\Big)
+\frac{2}{(2B+1)^{n}} \, ,
$$
we have
$$
\begin{aligned}
\AdvEUFSTATIC{\DSScheme^{\mathsf{d}}[\HF^{\mathsf{dlio}}]}{\advA}
\le\ &\AdvPRF{\PRF_0}{\advD_0}
+ q_s \cdot \frac{4\ell(\colsA+n)}{2B+1}
+ \sum_{j=1}^{q_s}\epsilon_j
+ \AdvLWE{q,n,\colsA,\ell,1}{\advD_{\mathsf{key}}} \\
&+ \AdvCSIS{\Fq_q,\chal_q,q,n,\colsA,\ell,\Bclip+B,\Bclip+B}{\advB_{\mathsf{main}}} \\
&+ \AdvIO{}{\advD^{\mathsf{io}}_{\mathsf{trig}}}
+ \sum_{i=3}^{4}\AdvPRF{\PRF_i}{\advD^{\mathsf{prf}}_{\mathsf{trig},i}}
+ \AdvLWE{q,n,\colsA,1,\Bclip}{\advD^{\mathsf{lwe}}_{\mathsf{trig}}} \\
&+ \alpha_{q,n}^{-1}\,\AdvSIS{q,n,\colsA+n+1,\Bclip}{\advB_{\mathsf{trig}}} \;.
\end{aligned}
$$
\end{theorem}

\paragraph*{Discussion.} Every assumption is polynomially hard and the scheme has no query parameter; the flooding condition still requires a quasi-polynomial modulus, as in Corollary~\ref{cor-dil-sis}. The two mechanisms are those of Theorem~\ref{thm-der-static}. The declared messages are known before key generation, so every puncture set is fixed in advance and nothing is guessed. And the trigger branch recomputes programmed points instead of storing them, so the published circuit has fixed size. Compared with Theorem~\ref{thm-der-static}, the lattice setting adds one LWE hop per direction of each per-query step, where Schnorr's output mask is exact, and the trigger-forgery case lands on SIS through an LWE hop rather than on DL.

\begin{proof}
Let $\advA_0$ declare $m_1, \ldots, m_{q_s}$ and $m^* \notin \{m_1,\ldots,m_{q_s}\}$; repeated declared messages are answered by replaying, and the hybrid runs over the distinct ones.

\heading{Game $\Gm_0$ (uniform nonces).} Replace each $\PRFEv_0(k_0, \enc{m_i,\kappa})$ by an independent uniform element of its range. The key $k_0$ appears nowhere else and the messages are declared before key generation, so one distinguisher bounds the change by $\AdvPRF{\PRF_0}{\advD_0}$. The game presamples all nonces before building the circuit.

\heading{Game $\Gm_1$ (rejection abort).} Abort, with output 0, if the first attempt of any signing query rejects; the cost is $q_s \cdot 4\ell(\colsA+n)/(2B+1)$ as in Game $\Gm_1$ of Theorem~\ref{thm-dil}. Each query is now answered by its first attempt, and its point $(\mat{w}_j, m_j)$ is a function of the presampled nonces alone.

\heading{Games $\Gm_{2,0}, \ldots, \Gm_{2,q_s}$ (trigger signing).} Game $\Gm_{2,i}$ answers the first $i$ distinct declared messages with the trigger record $(c_t(m_j), \mat{z}_t(m_j))$ recomputed from $k_3, k_4$, and the remaining ones honestly. Every game publishes the honest circuit $\CircD^{\mathsf{d}}$; hardwired circuits occur only inside the reductions between adjacent games. Game $\Gm_{2,0}$ is $\Gm_1$.

Fix $j$ and compare $\Gm_{2,j-1}$ with $\Gm_{2,j}$. Both endpoints add the event $\mathsf{Bad}_{\mathsf{eq}} := [\mat{w}_j = \mat{W}_{m_j}]$, with output 0 when it occurs; conditioned on everything else, $\mat{y}_{j,2}$ gives $\mat{w}_j$ min-entropy $n\log(2B+1)$ at any fixed target, so the two aborts cost at most $2(2B+1)^{-n}$, and intermediate games recompute the event from their current values. The chain is that of Appendix~\ref{app-der-static}, with the lattice substitutions. Replace the honest circuit by the two-point circuit that hardwires $\{(\enc{\mat{w}_j, m_j}''', v_j), (\enc{\mat{W}_{m_j}, m_j}''', c_t(m_j))\}$ and punctures $k_1, k_2$ at the two points and $k_3, k_4$ at $m_j$; the circuits compute the same function, because the trigger comparison at slot $m_j$ succeeds only at $\mat{W}_{m_j}$, which $\mathsf{Bad}_{\mathsf{eq}}$ separates from $\mat{w}_j$, and puncturing preserves functionality elsewhere, so iO bounds the change. Make the $\PRF_1$ and $\PRF_2$ values at the two points uniform, two puncturable-PRF hops with puncture sets fixed by the presampled nonces and the declared messages. Replace the masked point $\mat{w}_j + \mat{A}\mat{v}_1 + \mat{v}_2$ by uniform, one decision-LWE hop at box $B$ with one secret--noise pair, and then $v_j = \Fq_q(\text{uniform})$ is uniform on $\chal_q$ exactly, as in Game $\Gm_5$ of Theorem~\ref{thm-dil}. Re-derive the honest record by the exact rejection-sampling change of variables of Game $\Gm_6$ of Theorem~\ref{thm-dil}: with probability $1-\rho$ trigger the rejection abort, and otherwise sample $c_j \getsr \chal_q$ and $\mat{z}_j$ uniform on the clipped boxes and set $\mat{w}_j \gets \mat{A}\mat{z}_{j,1} + \mat{z}_{j,2} - \mat{T}c_j$. Make the $\PRF_3$ and $\PRF_4$ values at $m_j$ uniform, two more puncturable-PRF hops, so the trigger record is $(c_t, \mat{z}_t)$ uniform with $\mat{W}_{m_j}$ determined by the same identity. The two records are now independent and identically distributed, each a uniform challenge and a uniform clipped response with the point determined by the verification equation, conditioned on distinct points; the temporary circuit holds the canonically ordered set of the two point--value pairs and nothing that distinguishes them, so Lemma~\ref{lem-swap} applies with the verification identity in place of $R = g^{s}X^{-v}$, and the response to query $j$ exchanges. Reversing the steps with the second record in place of the first gives
$$
\big|\Pr[\Gm_{2,j-1} \Rightarrow 1] - \Pr[\Gm_{2,j} \Rightarrow 1]\big| \le \epsilon_j \;.
$$

\heading{Game $\Gm_3$ (uniform public key).} In $\Gm_{2,q_s}$ every signature is a trigger record and the secrets $\mat{S}_1, \mat{S}_2$ are unused. Replace $\mat{T}$ by uniform, at a cost of $\AdvLWE{q,n,\colsA,\ell,1}{\advD_{\mathsf{key}}}$.

\heading{A main-branch forgery.} Let $E_{\mathsf{main}}$ be the event that $\advA_1$ outputs a valid forgery $(c^*, \mat{z}^*)$ on $m^*$ whose hash value comes from the main branch. Adversary $\advB_{\mathsf{main}}$ embeds a $\CSIS$ instance as $(\mat{A}, \mat{T})$, simulates $\Gm_3$, and applies the extraction identity of Section~\ref{sec-dil-constr}, which gives a $\CSIS$ solution with norms at most $\Bclip + B$, nontrivial by the leading coordinate. Hence $\Pr[E_{\mathsf{main}}] \le \AdvCSIS{\Fq_q,\chal_q,q,n,\colsA,\ell,\Bclip+B,\Bclip+B}{\advB_{\mathsf{main}}}$.

\heading{A trigger-branch forgery.} Let $E_{\mathsf{trig}}$ be the event that the forgery's hash value comes from the trigger branch, so that $\mat{A}\mat{z}^*_1 + \mat{z}^*_2 - \mat{T}c^* = \mat{W}_{m^*}$ and $c^* = c_t(m^*)$. The slot $m^*$ is declared before key generation, so the reduction punctures $k_3, k_4$ at $m^*$ and publishes a circuit that, at slot $m^*$, compares against a hardwired point $\widehat{\mat{W}}$ and returns a hardwired $\widehat{c}$; with the true values this computes the same function, so iO bounds the switch, and two puncturable-PRF hops make $\widehat{c}$ uniform on $\chal_q$ and $\widehat{\mat{z}}$ uniform on the clipped boxes. One decision-LWE hop at box $\Bclip$ with one secret--noise pair replaces $\mat{A}\widehat{\mat{z}}_1 + \widehat{\mat{z}}_2$ by a uniform $\widehat{\mat{u}}$, so the hardwired point becomes $\widehat{\mat{W}} = \widehat{\mat{u}} - \mat{T}\widehat{c}$. Signing never needs slot $m^*$, since $m^*$ is not declared. A forgery in the trigger branch now satisfies $\mat{A}\mat{z}^*_1 + \mat{z}^*_2 = \widehat{\mat{u}}$ with both norms at most $\Bclip$. Adversary $\advB_{\mathsf{trig}}$ receives a SIS instance $[\mat{A}_1 \mid \mat{A}_2 \mid \mat{a}] \getsr \Zq^{n\times(\colsA+n+1)}$, outputs $\bot$ if $\mat{A}_2$ is not invertible, which happens with probability $1 - \alpha_{q,n}$, and otherwise sets $\mat{A} \gets \mat{A}_2^{-1}\mat{A}_1$ and $\widehat{\mat{u}} \gets \mat{A}_2^{-1}\mat{a}$, both uniform, and samples $\mat{T}$ and everything else itself. From the forgery it returns $(\mat{z}^*_1, \mat{z}^*_2, -1)$, which satisfies $\mat{A}_1\mat{z}^*_1 + \mat{A}_2\mat{z}^*_2 - \mat{a} = \mat{0}$, is nonzero, and has norm at most $\Bclip$. Hence
$$
\Pr[E_{\mathsf{trig}}] \le
\AdvIO{}{\advD^{\mathsf{io}}_{\mathsf{trig}}}
+ \sum_{i=3}^{4}\AdvPRF{\PRF_i}{\advD^{\mathsf{prf}}_{\mathsf{trig},i}}
+ \AdvLWE{q,n,\colsA,1,\Bclip}{\advD^{\mathsf{lwe}}_{\mathsf{trig}}}
+ \alpha_{q,n}^{-1}\,\AdvSIS{q,n,\colsA+n+1,\Bclip}{\advB_{\mathsf{trig}}} \;.
$$

\heading{Conclusion.} A valid forgery takes one of the two branches, and collecting the terms gives the stated bound.
\end{proof}

\adam{All: Five spots to check independently. (1) The min-entropy bound $(2B+1)^{-n}$ for $\mathsf{Bad}_{\mathsf{eq}}$: it conditions on everything independent of $\mat{y}_{j,2}$, and $\mat{W}_{m_j}$ is a function of the PRF keys, which are independent of $\mat{y}_{j,2}$; confirm no hop breaks this independence. (2) The change-of-variables step inside the per-query hop reuses the exactness argument of Game $\Gm_6$ of Theorem~\ref{thm-dil} for a single query; confirm the conditioning on the rejection abort composes with the per-query structure. (3) The trigger-forgery LWE hop is at box $\Bclip$, not $B$, because the trigger response lives in the clipped boxes; confirm the prelims LWE statement covers both boxes. (4) Puncture sets: $\enc{\mat{w}_j, m_j}'''$ depends on presampled nonces and $\mat{T}$; both exist before the hash key is built; confirm this meets Definition~\ref{def-pprf}. (5) The SIS embedding's invertibility loss duplicates the pattern of Lemma~\ref{lem-csis-sis}; confirm the norm bookkeeping $\max(\Bclip, 1) = \Bclip$.}

\subsection{Generic Conversion Functions and the Path to $B_\tau$} \label{sec-dil-generic}

The construction and theorem above fix the conversion function to $\Fq_q$. This subsection separates what is generic from what is not. The machinery of the scheme, the circuit, and the proof works for an abstract conversion function with three quantified properties. The gadget reduction of Lemma~\ref{lem-csis-sis} does not generalize, and for Dilithium's sparse challenge set it cannot, for an information-theoretic reason we state below. For that challenge set the circular assumption has to be assumed rather than reduced, and we discuss what a plausible conversion function looks like there.

\begin{definition} \label{def-convfun}
Let $q, n, \ell, \omega, \nu \in \N$ and $\beta \in [0,1]$. Let $\chal \subseteq \Z^{\ell}$ be a finite set with $\mat{0} \notin \chal$ and with $\|c\|_\infty \le \omega$ and $\|c\|_1 \le \nu$ for every $c \in \chal$. An efficient function $f \colon \Zq^n \to \chal$ is a \emph{$(\chal, \omega, \nu, \beta)$-conversion function} if the statistical distance between $f(\mat{w})$ for uniform $\mat{w} \in \Zq^n$ and the uniform distribution on $\chal$ is at most $\beta$.
\end{definition}

\paragraph*{Discussion.} The function $\Fq_q$ is a $(\chal_q, 1, \ell, 0)$-conversion function. Bijectivity onto $\chal_q$ is the case $\beta = 0$, the leading coordinate removes the zero vector, and binary coordinates bound both norms. The gadget identity $\mat{G}_q \Fq_q(\mat{w}) = \mat{w}$ is deliberately not part of the definition, because it is the one property that cannot generalize.

Each parameter has one job in the proofs of this part. The condition $\mat{0} \notin \chal$ is what lets verification skip a zero-challenge check and keeps every forgery nontrivial. The norms $\omega$ and $\nu$ set the size bounds: $\nu$ replaces $\ell$ in the response bound $\Bclip = B - \nu$, in the requirement $B \ge 4\nu(\colsA+n)$, in the shift bounds of Lemma~\ref{lem-dil-sim}, and in the final cSIS solution norms, which become $\Bclip + B$ as before. The parameter $\beta$ enters wherever the proofs use that a hash value is uniform on the challenge space: the masked branch of the circuit and the simulated challenge of Lemma~\ref{lem-dil-sim} each become $\beta$-close to uniform instead of exactly uniform, so the bounds of Theorem~\ref{thm-dil-bk} and Theorem~\ref{thm-dil} each gain a term of the form $(q_s+1)\beta$. Inspection shows the proofs use the four listed properties of $\Fq_q$ and nothing else, so the scheme, the hash construction, and both theorems generalize to any $(\chal, \omega, \nu, \beta)$-conversion function, with $\CSIS_{f}$ over the chosen $\chal$ as the hypothesis of the static theorem. \adam{All: I have re-traced where each property is used, but before we state the generic theorem formally someone should verify the claim independently. The check is a pass over Lemma~\ref{lem-dil-sim} and the games of Theorem~\ref{thm-dil} replacing ``uniform on $\chal_q$'' by ``$\beta$-close to uniform on $\chal$'' and $\ell$ by $\nu$ in every norm.}

\paragraph*{Sparse challenges.} Dilithium's challenge set is $B_\tau$, the set of vectors in $\{-1,0,1\}^{256}$ with exactly $\tau$ nonzero coordinates, with $\tau \in \{39, 49, 60\}$ across the ML-DSA parameter sets and $|B_{60}| \approx 2^{257}$. As a challenge space it has $\omega = 1$ and $\nu = \tau$, and $\tau$ is far smaller than $\ell$, so sparse challenges make every norm bound above smaller. This is exactly why Dilithium uses them. The zero vector is excluded by definition, since every element has $\tau$ nonzero coordinates.

No conversion function into $B_\tau$ admits a reduction in the style of Lemma~\ref{lem-csis-sis}. Such a reduction requires a matrix $\mat{G}_f$ with $\mat{G}_f f(\mat{w}) = \mat{w}$ for every $\mat{w} \in \Zq^n$, and the existence of such a matrix makes $f$ injective, so the challenge space must have at least $q^n$ elements. The set $B_{60}$ has about $2^{257}$ elements, while $q^n$ exceeds $2^{5000}$ already at Dilithium's dimensions. The obstruction is information-theoretic, and it does not depend on how $f$ is built.

So over $B_\tau$ the circular assumption stands on its own. This position is not new. The random-oracle analysis supporting Dilithium's standardization assumes SelfTargetMSIS~\cite{EC:KilLyuSch18}, a circular problem stated relative to the hash oracle, rather than reducing it to MSIS, and $\CSIS_f$ over $B_\tau$ with a concrete $f$ is the standard-model counterpart of that position. It is falsifiable and non-interactive. What it requires of $f$, beyond the definition above, is that no single challenge have a dense preimage: if some $c^*$ had $f^{-1}(c^*)$ of noticeable density, an attacker could sample a short $\mat{z}$, set $\mat{w} \gets \mat{A}\mat{z}_1 + \mat{z}_2 - \mat{T}c^*$, and test $f(\mat{w}) = c^*$, and success is a cSIS solution. A $\beta$-regular $f$ with $\beta$ negligible has every preimage of density about $1/|B_\tau| \approx 2^{-257}$, which closes this attack. Candidate conversion functions, in decreasing order of explicitness: Dilithium's own $\mathsf{SampleInBall}$~\cite{Dilithium} applied to low-order bits of the coordinates of $\mat{w}$, which is unkeyed and fully explicit, and for which the seed bits of a uniform $\mat{w}$ are within $n/2q$ of uniform; $\mathsf{SampleInBall}$ composed with a concrete hash function, which is the deployed shape; and a keyed $f$, for which the assumption becomes average-case over the key. Stating $\CSIS_f$ over $B_\tau$ for the first candidate and collecting evidence for it is the main open problem of this part.

\paragraph*{Modules and compression.} Over $R_q = \Z_q[X]/(X^{256}+1)$ the template becomes cMSIS, with the matrices over $R_q$ and shortness read coefficientwise. Decomposition challenges do not transfer well to this setting. The lifted gadget argument is sound, with the invertibility constant of Lemma~\ref{lem-csis-sis} recomputed over $R_q$, but ring multiplication makes the response bounds scale with the $\ell_1$ norm of the challenge counted over every coefficient of every ring entry. A decomposition challenge has $\mathrm{rank}\cdot\lceil\log_2 q\rceil$ ring entries of $256$ coefficients each, more than twenty thousand nonzero coefficients at ML-DSA dimensions, against $\nu = \tau \le 60$ for $B_\tau$. The parameter advantage that motivates modules is spent entirely on the challenge, so we do not pursue provable cMSIS over decomposition challenges in the module setting. The module case that matters is $B_\tau$ under assumed cMSIS with a concrete conversion function, and there the obstruction above is unchanged by the module structure. Compression of the commitment is a separate gap that no choice of conversion function affects: Dilithium computes the challenge from $\mathsf{HighBits}(\mat{w})$ and verification reconstructs it through a hint vector, so the extraction identity holds only up to a rounding term, and our proofs use the exact identity. Section~\ref{sec-conclusion} records both.

\paragraph*{Outlook: the adaptive port.} Section~\ref{sec-dil-derand} completes the static half of the derandomized port. What remains from Part~I is the adaptive result. Pre-hash the message with the tunable lossy function as in Section~\ref{sec-derand} and guess digests from the enumerated lossy image; the guessing combinatorics do not touch the algebra and should transfer. The lattice-specific work is the interaction of the digest guessing with the flooding condition and with the LWE hops at the trigger points, and circuit equivalence on malformed inputs. The pre-hash is a second, independent gap. The proof guesses the digest of a query from the enumerated lossy image, so it needs a tunable lossy function whose lossy image can be enumerated, and the only instantiation we have is Zhandry's ELF cascade, whose security is a DDH-type assumption and is broken by quantum attacks. Completing the port with it would give a standard-model adaptive theorem for the derandomized lattice scheme whose assumption profile is not post-quantum, and this caveat differs in kind from the iO caveat this part already has, since iO has plausible lattice-based candidates while exponential DDH has no post-quantum version. A post-quantum lossy function with enumerable image is the open problem that removes the gap, and Section~\ref{sec-conclusion} records it. Deterministic ML-DSA is the signing mode such a theorem would cover; the hedged default mode stays open, as it does for BIP340 on the Schnorr side. A proof would have to establish the trigger response distribution in the clipped box, its exchange with an honest rejection-sampled response, the treatment of rejected attempts, the interaction with the flooding condition or a replacement for it, the LWE hops at the trigger points, and circuit equivalence on malformed inputs. We state this as a program, not as a theorem whose proof is routine. The trigger comparison $R = g^{s_t}X^{-c_t}$ becomes $\mat{w} = \mat{A}\mat{z}_{t,1} + \mat{z}_{t,2} - \mat{T}c_t$ with the PRF-derived response sampled in the clipped box. \adam{All: Beyond the ports: compression and hints, the sparse challenge set $B_\tau$ (the working draft's open problem), balanced decompositions at prime $q$, the module lift, and derandomized signing. Suggested split: Ojaswi+Carter on the key-recovery port and parameters, Xinyu on the sparse-challenge assumption ($\CSIS$ over $B_\tau$: no gadget-style reduction can exist, since a challenge in $B_\tau$ cannot determine $\mat{w}$; the question is a plausible explicit $f$ and evidence), Alex+Adam on the derandomized ports and the module lift.}
